\subsection{Đánh giá và thảo luận}
\label{sec:dsom_discussion}

\subsubsection{Nhận xét về phương pháp}

DSoM khai thác thông tin
đa phương thức trong huấn luyện teacher nhưng student chỉ dùng cấu trúc
khi suy luận. Còn RC lựa chọn nguồn hướng dẫn cho từng mẫu, trong khi
NDKD phân tách thông tin được truyền từ teacher.

Ở mục~5.5 của bài báo~\cite{10.1145/3746027.3754576}, tác giả
phân tích tỷ lệ lựa chọn các tổ hợp teacher. Theo nhóm, các
thống kê này giúp mô tả hành vi của policy, nhưng sự đa dạng
trong lựa chọn chưa đủ chứng minh agent đã loại bỏ được thông
tin nhiễu. Cần đối chiếu lựa chọn của agent với chất lượng dự
đoán của teacher và kết quả student trên cùng nhóm truy vấn.

\subsubsection{Phạm vi kiểm chứng của kết quả tái lập}

Các so sánh ở Bảng~\ref{tab:dsom_table4} và
Bảng~\ref{tab:dsom_ablation} chưa cho thấy learned RC có ưu
thế ổn định so với những lựa chọn đơn giản trong các lần chạy
của nhóm. Nhận xét này áp dụng cho những cấu hình đã chạy
với một random seed. Nhóm chưa thể tách riêng nguyên nhân
của các chênh lệch vì quá trình tái lập còn có những sai khác
về dữ liệu, đặc trưng và cách triển khai.

Nhóm cũng phân biệt việc chạy lại mã nguồn với việc kiểm
chứng đầy đủ mô tả phương pháp. Cấu hình \texttt{source\_frozen}
không kiểm chứng được quá trình học policy của tác giả;
sai khác giữa trung bình và tổng trong NEKD còn làm thay đổi
mục tiêu distillation. Vì vậy, metric gần với paper chưa đủ
xác nhận hai quy trình tương đương. Các giới hạn về validation
DB15K, đặc trưng ảnh tái tạo và giao thức đánh giá đã nêu
cần được xem xét khi đối chiếu kết quả.